
\section{考察}
% ここに結果の解釈や誤差要因、改善点を記載します。
% 実験結果を理論と比較し、差異があればその原因を考察します。

\subsection{結果の解釈}
% 測定結果が何を意味するのかを説明します。
% 理論値との比較を行い、一致・不一致の理由を考察します。

\subsection{誤差要因}
% 測定誤差の原因を具体的に挙げます。
% 例: 測定機器の精度、環境条件、操作上の誤差など

\subsection{改善案}
% 実験の精度を向上させるための改善案を提案します。
